\documentclass{article}
\usepackage[top=2cm,bottom=2cm,left=3cm,right=3cm]{geometry}
\usepackage{enumitem}
\usepackage{xcolor}
\usepackage{minted}
\usepackage[most]{tcolorbox}
\tcbuselibrary{minted,breakable}
\usepackage{fontspec}
\usepackage{polyglossia}
\setdefaultlanguage{ukrainian}
\setotherlanguages{english}
\usepackage{Alegreya}
\usepackage{FiraMono}
\newfontfamily\cyrillicfont{Alegreya}
\newfontfamily\cyrillicfontsf{Alegreya Sans}
\newfontfamily\cyrillicfonttt{DejaVu Sans Mono}
\usepackage{sciffi-python}
\usepackage{sciffi-memo}

\usemintedstyle{emacs}

\newtcolorbox{CodeBox}[1][]{
    breakable,
    enhanced,
    colback=gray!10,
    colframe=gray!70,
    fonttitle=\bfseries,
    title=#1,
    sharp corners,
    boxrule=0.8pt,
    top=2mm,
    bottom=2mm,
    left=2mm,
    right=2mm,
    fontupper=\footnotesize,
}

\newtcolorbox{ResultBox}[1][]{
    breakable,
    enhanced,
    colback=gray!5,
    colframe=gray!50,
    fonttitle=\bfseries,
    title=#1,
    sharp corners,
    boxrule=0.6pt,
    top=1mm,
    bottom=1mm,
    left=2mm,
    right=2mm,
    fontupper=\footnotesize,
}

\begin{document}
    \begin{titlepage}
        \begin{center}
            {\Large
                МІНІСТЕРСТВО ОСВІТИ ТА НАУКИ УКРАЇНИ

                НАЦІОНАЛЬНИЙ ТЕХНІЧНИЙ УНІВЕРСИТЕТ УКРАЇНИ
                «КИЇВСЬКИЙ ПОЛІТЕХНІЧНИЙ ІНСТИТУТ ІМЕНІ ІГОРЯ СІКОРСЬКОГО»

                ФАКУЛЬТЕТ ІНФОРМАТИКИ ТА ОБЧИСЛЮВАЛЬНОЇ ТЕХНІКИ

                КАФЕДРА ІНФОРМАЦІЙНИХ СИСТЕМ ТА ТЕХНОЛОГІЙ
            }


            \vspace{60mm}
            {\large
                \textbf{ЗВІТ}

                \vspace{5mm}

                \textbf{До лабораторної роботи 3}

                \vspace{5mm}

                з дисципліни «Технології Розробки Програмного Забезпечення»
            }

            На тему «Музичний Магазин»

            варіант \textnumero 14 (\textnumero 4)
        \end{center}

        \vfill
        \hfill
        \begin{minipage}[t]{0.35\textwidth}
            \textbf{Виконав:}
        \end{minipage}

        \vfill

        \begin{center}
            {\bf
                Київ

                КПІ Ім. Ігоря Сікорського

                2025
            }
        \end{center}
    \end{titlepage}
    \newpage

    Тема: з’єднання та теоретико-множинні операції над відношенням.

    Мета: вивчити команди формування запитів до кількох таблиць бази даних (JOIN).

    Навички що будуть здобуті: вміння формувати такі запити до бази даних, що використовують з’єднання таблиць неявним та явним чином.

    \section*{Хід виконання:}
    \begin{sciffi}{python}
        from pathlib import Path
        from subprocess import run, PIPE
        from collections.abc import Sequence
        from contextlib import contextmanager
        
        @contextmanager
        def env(name: str, *args: str):
            print(f"\\begin{{{name}}}" + "".join(args))
            try:
                yield
            finally:
                print(f"\\end{{{name}}}")

        def sh(*args: str | Path) -> str:
            return run(
                tuple(map(str, args)),
                stdout=PIPE,
                stderr=PIPE,
                text=True,
            ).stdout

        scripts = Path("./scripts/")
        # Find all .sql files and sort them numerically by their filename
        sql_files = sorted(scripts.glob("*.sql"), key=lambda p: int(p.stem))

        flow: Sequence[tuple[Path, Sequence[str]]] = [
            (script, ("--no-echo",)) for script in sql_files
        ]

        for script, args in flow:
            source = script.read_text(encoding="utf-8")
            output = sh("tsqlexec", *args, script)

            print(f"\\subsection*{{Запит з файлу: {script.name}}}")
            with env("CodeBox", "[SQL Source]"):
                with env("minted", "[breaklines,linenos]", "{sql}"):
                    print(source, end="")
            print()

            if not output.strip():
                continue

            print("Result:")
            with env("ResultBox", "[Output]"):
                with env("minted", "[breaklines]", "{text}"):
                    print(output, end="")
            print("\\vspace{5mm}")

    \end{sciffi}

    \newpage
    \section*{Контрольні питання}
    \begin{enumerate}[label=\arabic*., leftmargin=*]
        \item \textbf{Поняття операції вибірки.}
        
        Операція вибірки (Selection) — це одна з фундаментальних операцій реляційної алгебри. Вона дозволяє відфільтрувати рядки (кортежі) з таблиці, які задовольняють певній умові (предикату). В мові SQL ця операція реалізується за допомогою речення \texttt{WHERE}.
        \begin{CodeBox}[Приклад вибірки]
        \begin{minted}[breaklines]{sql}
-- Вибрати всіх клієнтів з міста Київ
SELECT * FROM Customers WHERE City = 'Київ';
        \end{minted}
        \end{CodeBox}

        \item \textbf{Поняття операції проекції.}
        
        Операція проекції (Projection) — це операція реляційної алгебри, яка вибирає певні стовпці (атрибути) з таблиці, відкидаючи інші. В SQL вона реалізується переліком стовпців у частині \texttt{SELECT}.
        \begin{CodeBox}[Приклад проекції]
        \begin{minted}[breaklines]{sql}
-- Вибрати лише імена та прізвища клієнтів
SELECT FirstName, LastName FROM Customers;
        \end{minted}
        \end{CodeBox}
        
        \item \textbf{Поняття та синтаксис операції Декартів добуток.}
        
        Декартів добуток (Cartesian Product) — це операція, яка комбінує кожен рядок з однієї таблиці з кожним рядком з іншої таблиці. Результатом є таблиця, що містить усі можливі пари рядків. В SQL декартів добуток утворюється, коли таблиці перелічуються у \texttt{FROM} без умови з'єднання.
        \begin{CodeBox}[Синтаксис Декартового добутку]
        \begin{minted}[breaklines]{sql}
SELECT * FROM Table1, Table2;
-- Або з використанням CROSS JOIN
SELECT * FROM Table1 CROSS JOIN Table2;
        \end{minted}
        \end{CodeBox}
        Зазвичай це призводить до дуже великих результуючих наборів і використовується рідко, частіше за все — як основа для подальшої фільтрації.

        \item \textbf{Властивості та застосування псевдонімів.}
        
        Псевдоніми (aliases) — це тимчасові імена, що надаються таблицям або стовпцям у межах одного запиту. Їх використовують для:
        \begin{itemize}
            \item \textbf{Скорочення довгих імен:} \texttt{FROM VeryLongTableName AS vlt}.
            \item \textbf{Покращення читабельності:} \texttt{SELECT p.Name, c.CategoryName FROM Products AS p JOIN Categories AS c ON ...}.
            \item \textbf{Уникнення неоднозначності:} коли з'єднуються таблиці з однаковими іменами стовпців (\texttt{Orders.ID}, \texttt{Customers.ID}).
            \item \textbf{Самоз'єднання (Self-Join):} коли таблиця з'єднується сама з собою, псевдоніми є обов'язковими для розрізнення екземплярів таблиці.
        \end{itemize}
        
        \item \textbf{Явний та неявний синтаксис.}
        
        Це стосується синтаксису з'єднання таблиць:
        \begin{itemize}
            \item \textbf{Неявний синтаксис (старий стиль):} Таблиці перелічуються у реченні \texttt{FROM} через кому, а умови з'єднання вказуються у \texttt{WHERE}.
            \begin{CodeBox}[Неявний синтаксис]
            \begin{minted}{sql}
SELECT c.Name, o.OrderDate FROM Customers c, Orders o WHERE c.ID = o.CustomerID;
            \end{minted}
            \end{CodeBox}
            \item \textbf{Явний синтаксис (стандарт ANSI SQL-92):} Використовується ключове слово \texttt{JOIN}, а умови з'єднання вказуються у реченні \texttt{ON}.
            \begin{CodeBox}[Явний синтаксис]
            \begin{minted}{sql}
SELECT c.Name, o.OrderDate FROM Customers c JOIN Orders o ON c.ID = o.CustomerID;
            \end{minted}
            \end{CodeBox}
        \end{itemize}
        Явний синтаксис є кращим, оскільки він чітко відокремлює логіку з'єднання від логіки фільтрації.

        \item \textbf{Синтаксис команди JOIN в пропозиції FROM.}
        
        Загальний синтаксис команди \texttt{JOIN} виглядає так:
        \begin{CodeBox}[Загальний синтаксис JOIN]
        \begin{minted}{sql}
FROM table1
[join_type] JOIN table2 ON join_condition
...
        \end{minted}
        \end{CodeBox}
        Де \texttt{join\_type} — це тип з'єднання (\texttt{INNER}, \texttt{LEFT}, \texttt{RIGHT}, \texttt{FULL}), а \texttt{join\_condition} — умова, за якою рядки з'єднуються.

        \item \textbf{Які бувають види з’єднань.}
        
        Основні види з'єднань:
        \begin{itemize}
            \item \textbf{INNER JOIN (Внутрішнє з'єднання):} Повертає тільки ті рядки, для яких знайдено збіг в обох таблицях.
            \item \textbf{LEFT JOIN (Ліве зовнішнє з'єднання):} Повертає всі рядки з лівої таблиці та відповідні рядки з правої. Якщо збігу немає, стовпці правої таблиці заповнюються \texttt{NULL}.
            \item \textbf{RIGHT JOIN (Праве зовнішнє з'єднання):} Повертає всі рядки з правої таблиці та відповідні рядки з лівої. Якщо збігу немає, стовпці лівої таблиці заповнюються \texttt{NULL}.
            \item \textbf{FULL OUTER JOIN (Повне зовнішнє з'єднання):} Повертає всі рядки з обох таблиць. Якщо збігу немає, відповідні стовпці заповнюються \texttt{NULL}.
            \item \textbf{CROSS JOIN (Перехресне з'єднання):} Повертає декартів добуток таблиць.
        \end{itemize}

        \item \textbf{Поняття та синтаксис внутрішнього з’єднання.}
        
        \textbf{Поняття:} \texttt{INNER JOIN} повертає рядки, коли є збіг значень у з'єднуваних стовпцях обох таблиць. Це найпоширеніший тип з'єднання.
        
        \textbf{Синтаксис:}
        \begin{CodeBox}[Синтаксис INNER JOIN]
        \begin{minted}{sql}
SELECT columns FROM table1 INNER JOIN table2 ON table1.column = table2.column;
        \end{minted}
        \end{CodeBox}
        
        \item \textbf{Поняття та синтаксис лівого зовнішнього з’єднання.}
        
        \textbf{Поняття:} \texttt{LEFT OUTER JOIN} повертає всі рядки з лівої таблиці та збіги з правої. Якщо для рядка з лівої таблиці не знайдено збігу в правій, результат для правої сторони буде містити \texttt{NULL}.
        
        \textbf{Синтаксис:}
        \begin{CodeBox}[Синтаксис LEFT JOIN]
        \begin{minted}{sql}
SELECT columns FROM table1 LEFT JOIN table2 ON table1.column = table2.column;
        \end{minted}
        \end{CodeBox}

        \item \textbf{Поняття та синтаксис правого зовнішнього з’єднання.}
        
        \textbf{Поняття:} \texttt{RIGHT OUTER JOIN} є протилежністю лівого з'єднання. Воно повертає всі рядки з правої таблиці та збіги з лівої. Якщо для рядка з правої таблиці не знайдено збігу, результат для лівої сторони буде містити \texttt{NULL}.
        
        \textbf{Синтаксис:}
        \begin{CodeBox}[Синтаксис RIGHT JOIN]
        \begin{minted}{sql}
SELECT columns FROM table1 RIGHT JOIN table2 ON table1.column = table2.column;
        \end{minted}
        \end{CodeBox}

        \item \textbf{Поняття та синтаксис повного зовнішнього з’єднання.}
        
        \textbf{Поняття:} \texttt{FULL OUTER JOIN} об'єднує результати \texttt{LEFT} та \texttt{RIGHT JOIN}. Воно повертає всі рядки з обох таблиць. Якщо для рядка з однієї таблиці немає збігу в іншій, відповідні стовпці заповнюються \texttt{NULL}.
        
        \textbf{Синтаксис:}
        \begin{CodeBox}[Синтаксис FULL OUTER JOIN]
        \begin{minted}{sql}
SELECT columns FROM table1 FULL OUTER JOIN table2 ON table1.column = table2.column;
        \end{minted}
        \end{CodeBox}
        
    \end{enumerate}
\end{document}
